\documentclass[../../main.tex]{subfiles}
\begin{document}

\begin{center}
  {\fontsize{48pt}{58pt}\selectfont \textbf{1971}\normalsize\textbf{年度}}
\end{center}

\begin{table}[htb]
  \centering
  \caption{出題テーマと難易度}
  \begin{tabular}{|c|p{6.2cm}|c|c|c|}
    \hline
    \textbf{問題番号} & \textbf{テーマ} & \textbf{思考量} & \textbf{計算量} & \textbf{総合} \\
    \hline
    第1問 & 媒介変数で表された点と原点の距離 & 易 & 普 & 易 \\
    \hline
    第2問 & 平方を用いた漸化式と単調性 & 難 & 普 & 難 \\
    \hline
    第3問 & 積分の二乗誤差を最小にする一次式 & 易 & 普 & 普 \\
    \hline
    第4問 & 指数関数の部分和と実根の個数 & 易 & 普 & 易 \\
    \hline
    第5問 & 円周上の等分点までの距離の二乗平均 & 易 & 易 & 易 \\
    \hline
    第6問 & 三人のじゃんけんと勝者決定の確率 & 易 & 易 & 易 \\
    \hline
  \end{tabular}
\end{table}

\begin{table}[htb]
  \centering
  \caption{解答}
  \renewcommand{\arraystretch}{1.4}
  \begin{tabular}{|c|c|p{9cm}|}
    \hline
    \textbf{問題番号} & \textbf{小問} & \textbf{答え} \\
    \hline
    第1問 & -- & 最大値は$2$（$t=\dfrac{\pi}{6}$），最小値は$1$（$t=\dfrac{2\pi}{3}$） \\
    \hline
    第2問 & (1) & $x\ne2$なら数列は狭義単調増加（証明問題） \\
    \cline{2-3}
    & (2) & $x<2$なら$a_n<1$，かつ$2-x>n\epsilon^2$（証明問題） \\
    \hline
    第3問 & -- & $(a,b)=(24,-10)$ \\
    \hline
    第4問 & -- & $f_n'=f_{n-1}$．$n$が奇数なら実根は$1$個，偶数なら$0$個（証明問題） \\
    \hline
    第5問 & -- & $\lambda_P=1+a^2+b^2$．$P=(0,0)$で最小 \\
    \hline
    第6問 & -- & $(2k-1)\left(\dfrac{1}{3}\right)^k$ \\
    \hline
  \end{tabular}
\end{table}

\paragraph{試験形態}
6問．試験時間は要確認．

\end{document}
